\documentclass{beamer}

\usepackage[T2A]{fontenc}
\usepackage[utf8]{inputenc}
\usepackage[english,russian]{babel}
\input{settings.tex}


\title{$H_\infty$, Юнг}
\subtitle{Вычислительный интеллект, Лекция 10}
\author{Сергей Савин}
\centering
\date{\mydate}



\begin{document}
\maketitle


\begin{frame}{\texorpdfstring{$\mathcal{H}_\infty$}{} управление, функции накопления и KYP}
	\begin{flushleft}
		
		{\small
			\textbf{Система:}
			\[
			\dot{x} = Ax + Bw, \quad y = Cx + Dw
			\]
			
			
			\textbf{$\mathcal{H}_\infty$ критерий:} обеспечить
			\[
			\|y\|_2 \le \gamma \|w\|_2 \quad \text{для всех } w \in \mathcal{L}_2
			\]
			
			
			\textbf{Функция накопления (неравенство диссипации):} \\
			Найти $V(x) = x^\top P x$, $P \succ 0$, такую что
			\[
			\dot{V}(x) + y^\top y - \gamma^2 w^\top w \le 0
			\]
			
			\textbf{KYP (Лемма об ограниченной вещественности):} \\
			Вышеуказанное выполняется тогда и только тогда, когда $\exists P \succ 0$ такое что
			\[
			\begin{bmatrix}
				A^\top P + PA & PB & C^\top \\
				B^\top P & -\gamma^2 I & D^\top \\
				C & D & -I
			\end{bmatrix} \prec 0
			\]
			
			
			\textbf{Вывод:} $\Rightarrow \|G(s)\|_\infty < \gamma$}
		
	\end{flushleft}
	
\end{frame}


\begin{frame}{H бесконечность}
	\begin{flushleft}
		
		Рассмотрим неопределённую систему:
		\begin{align*}
			\dot{x} &= (A + \Delta A)x + Bu + Dw, \\
			y &= Cx,
		\end{align*}
		со структурированной неопределённостью
		\[
		\Delta A = N \Delta M, \quad \|\Delta\| < 1.
		\]
		
		Применим обратную связь по состоянию:
		\[
		u = Kx,
		\]
		что приводит к замкнутой системе:
		\[
		\dot{x} = (A + BK + \Delta A)x + Dw.
		\]
		
	\end{flushleft}
\end{frame}



\begin{frame}{Шаг 1: Лемма KYP}
	\begin{flushleft}
		
		Условие $\mathcal{H}_\infty$ выполняется, если существует $P \succ 0$ такое что:
		\[
		\begin{bmatrix}
			(A + BK + \Delta A)^\top P + P(A + BK + \Delta A) & PD & C^\top \\
			D^\top P & -\gamma I & 0 \\
			C & 0 & -\gamma I
		\end{bmatrix} \prec 0.
		\]
		
		
	\end{flushleft}
\end{frame}


\begin{frame}{Шаг 2: Разделение номинальной и неопределённой частей}
	\begin{flushleft}
		
		
		Обозначим:
		\[
		W_1 =
		\begin{bmatrix}
			(A + BK)^\top P + P(A + BK) & PD & C^\top \\
			D^\top P & -\gamma I & 0 \\
			C & 0 & -\gamma I
		\end{bmatrix},
		\]
		\[
		W_2 =
		\begin{bmatrix}
			(\Delta A)^\top P + P(\Delta A) & 0 & 0 \\
			0 & 0 & 0 \\
			0 & 0 & 0
		\end{bmatrix}.
		\]
		
		Используя $\Delta A = N \Delta M$, получаем:
		\[
		W_2 =
		\begin{bmatrix}
			M^\top \Delta^\top N^\top P + PN \Delta M & 0 & 0 \\
			0 & 0 & 0 \\
			0 & 0 & 0
		\end{bmatrix}.
		\]
		
	\end{flushleft}
\end{frame}


\begin{frame}{Неравенство Юнга (матричная форма)}
	
	\begin{flushleft}
		
		\textbf{Формулировка:} Для любых матриц $X, Y$ согласованных размерностей и любого $\varepsilon > 0$,
		\[
		X^\top Y + Y^\top X \;\preceq\; \varepsilon X^\top X + \varepsilon^{-1} Y^\top Y
		\]
		
		
		\medskip
		
		\textbf{Эквивалентная форма (используется в робастном управлении):}
		\[
		X \Delta Y + (X \Delta Y)^\top
		\;\preceq\;
		\varepsilon XX^\top + \varepsilon^{-1} Y^\top Y,
		\quad \|\Delta\| \le 1
		\]
		
		\medskip
		
		\textbf{Роль в LMI:}
		\begin{itemize}
			\item Заменяет неопределённые члены \emph{выпуклыми оценками}
			\item Вводит скалярный множитель $\varepsilon$
		\end{itemize}
		
	\end{flushleft}
	
\end{frame}



\begin{frame}{Факторизация неопределённого члена}
	\begin{flushleft}
		
		
		
		\textbf{Неопределённый блок:}
		\[
		W_2 =
		\begin{bmatrix}
			M^\top \Delta^\top N^\top P + PN \Delta M & 0 & 0 \\
			0 & 0 & 0 \\
			0 & 0 & 0
		\end{bmatrix}
		\]
		
		\medskip
		
		\textbf{Определим:}
		\[
		V_1 =
		\begin{bmatrix}
			PN \\ 0 \\ 0
		\end{bmatrix},
		\quad
		V_2 =
		\begin{bmatrix}
			M^\top \\ 0 \\ 0
		\end{bmatrix}
		\]
		
		\medskip
		
		\textbf{Компактная факторизация:}
		\[
		W_2 = V_1 \Delta V_2^\top + V_2 \Delta^\top V_1^\top
		\]
		
		\medskip
		
		\textbf{Развёрнутая форма:}
		\[
		W_2 =
		\begin{bmatrix}
			PN \\ 0 \\ 0
		\end{bmatrix}
		\Delta
		\begin{bmatrix}
			M & 0 & 0
		\end{bmatrix}
		+
		\begin{bmatrix}
			M^\top \\ 0 \\ 0
		\end{bmatrix}
		\Delta^\top
		\begin{bmatrix}
			N^\top P & 0 & 0
		\end{bmatrix}
		\]
		
	\end{flushleft}
	
\end{frame}

\begin{frame}{Шаг 3: Неравенство Юнга}
	\begin{flushleft}
		
		
		Используя неравенство Юнга:
		\[
		W_2 = V_1 \Delta V_2^\top + V_2 \Delta^\top V_1^\top \preceq \frac{1}{\varepsilon} V_1 V_1^\top + \varepsilon V_2 V_2^\top,
		\]
		заменяем условие на:
		\[
		W_1 + \frac{1}{\varepsilon} V_1 V_1^\top + \varepsilon V_2 V_2^\top \prec 0.
		\]
		
		
	\end{flushleft}
\end{frame}


\begin{frame}{Шаг 4: Первое дополнение Шура}
	\begin{flushleft}
		
		Применяя дополнение Шура к $\frac{1}{\varepsilon} V_1 V_1^\top$:
		\[
		\begin{bmatrix}
			W_1 + \varepsilon V_2 V_2^\top & V_1 \\
			V_1^\top & -\varepsilon I
		\end{bmatrix}
		\prec 0.
		\]
		
		Раскрывая:
		\[
		\begin{bmatrix}
			(A + BK)^\top P + P(A + BK) + \varepsilon M^\top M & PD & C^\top & PN \\
			D^\top P & -\gamma I & 0 & 0 \\
			C & 0 & -\gamma I & 0 \\
			N^\top P & 0 & 0 & -\varepsilon I
		\end{bmatrix}
		\prec 0.
		\]
		
	\end{flushleft}
\end{frame}


\begin{frame}{Шаг 5: Замена переменных}
	\begin{flushleft}
		
		Пусть:
		\[
		Q = P^{-1}, \quad U = KQ.
		\]
		
		Применим конгруэнтное преобразование $\mathrm{diag}(Q, I, I, I)$:
		\[
		\begin{bmatrix}
			QA^\top + AQ + BU + U^\top B^\top + \varepsilon Q M^\top M Q
			& D & QC^\top & N \\
			D^\top & -\gamma I & 0 & 0 \\
			CQ & 0 & -\gamma I & 0 \\
			N^\top & 0 & 0 & -\varepsilon I
		\end{bmatrix}
		\prec 0.
		\]
		
		
	\end{flushleft}
\end{frame}



\begin{frame}{Шаг 6: Второе дополнение Шура}
	\begin{flushleft}
		
		\[
		\varepsilon Q M^\top M Q
		=
		(Q M^\top)(\varepsilon I)(M Q).
		\]
		
		Теперь применим дополнение Шура:
		\[
		\begin{bmatrix}
			QA^\top + AQ + BU + U^\top B^\top
			& D
			& QC^\top
			& N
			& Q M^\top \\
			D^\top
			& -\gamma I
			& 0
			& 0
			& 0 \\
			CQ
			& 0
			& -\gamma I
			& 0
			& 0 \\
			N^\top
			& 0
			& 0
			& -\varepsilon I
			& 0 \\
			M Q
			& 0
			& 0
			& 0
			& -\varepsilon^{-1} I
		\end{bmatrix}
		\prec 0.
		\]
		
	\end{flushleft}
\end{frame}



\begin{frame}{Шаг 7: Замена переменных}
	\begin{flushleft}
		
		Определим:
		\[
		\alpha = \varepsilon^{-1}.
		\]
		
		Затем умножим LMI на конгруэнтное преобразование:
		\[
		\mathrm{diag}(I, I, I, \alpha I, I),
		\]
		чтобы получить:
		
		\[
		\begin{bmatrix}
			QA^\top + AQ + BU + U^\top B^\top
			& D
			& QC^\top
			& \alpha N
			& Q M^\top \\
			D^\top
			& -\gamma I
			& 0
			& 0
			& 0 \\
			CQ
			& 0
			& -\gamma I
			& 0
			& 0 \\
			\alpha N^\top
			& 0
			& 0
			& -\alpha I
			& 0 \\
			M Q
			& 0
			& 0
			& 0
			& -\alpha I
		\end{bmatrix}
		\prec 0.
		\]
		
	\end{flushleft}
\end{frame}


clusion}
\begin{flushleft}

{\small 
	The resulting SDP:
	
	
	\begin{equation}
		\label{robust_stability_sdp}
		\hspace*{-0.5cm} 
		\begin{aligned}
			& \underset{Q, U, \gamma, \alpha}{\text{minimize}} 
			& &  0, \\
			& \text{subject to:} & & 
			\begin{cases}
				\begin{bmatrix}
					QA^\top + AQ + BU + U^\top B^\top
					& D
					& QC^\top
					& \alpha N
					& Q M^\top \\
					D^\top
					& -\gamma I
					& 0
					& 0
					& 0 \\
					CQ
					& 0
					& -\gamma I
					& 0
					& 0 \\
					\alpha N^\top
					& 0
					& 0
					& -\alpha I
					& 0 \\
					M Q
					& 0
					& 0
					& 0
					& -\alpha I
				\end{bmatrix}
				\prec 0, \\ 
				Q \succ 0, 
				\gamma \geq 0, \alpha \geq 0
			\end{cases}
		\end{aligned}
	\end{equation}
	
	
	with controller recovered as:
	\[
	K = U Q^{-1}.
	\]
	
}
\end{flushleft}
\end{frame}


\myqrframe


\end{document}
